\section{Appendix A: Full Prompt Template}

The following prompt was used verbatim for all eight models. The \texttt{\{step\}} placeholder was replaced with either ``Step 1'' or ``Step 2 CK''; the \texttt{\{question\}} and \texttt{\{options\}} placeholders were filled from the dataset.

\begin{verbatim}
You are taking the USMLE {step} examination.

Question:
{question}

Options:
A. {option_A}
B. {option_B}
C. {option_C}
D. {option_D}

Instructions:
- Think through this carefully, step by step.
- Consider all options before deciding.
- State your final answer as: ANSWER: [A/B/C/D]
- The ANSWER line must be the last thing you write.
\end{verbatim}

Temperature was set to 0 for all calls. No system prompt was used. No few-shot examples were provided.

\section{Appendix B: Subject Distribution of Sampled Questions}

\begin{table}[htbp]
\centering
\caption{Subject Distribution in Stratified Sample}
\label{tab:subject_distribution}
\begin{tabular}{lcc}
\toprule
Subject & Step 1 (N=100) & Step 2 CK (N=100) \\
\midrule
Anatomy & -- & -- \\
Physiology & -- & -- \\
Biochemistry & -- & -- \\
Pharmacology & -- & -- \\
Pathology & -- & -- \\
Microbiology & -- & -- \\
Immunology & -- & -- \\
Behavioral Science & -- & -- \\
Biostatistics & -- & -- \\
Internal Medicine & -- & -- \\
Surgery & -- & -- \\
Pediatrics & -- & -- \\
Obstetrics \& Gynecology & -- & -- \\
Psychiatry & -- & -- \\
Neurology & -- & -- \\
\midrule
Total & 100 & 100 \\
\bottomrule
\end{tabular}
\end{table}

% PLACEHOLDER: Replace -- values with actual counts after running 02_sample_questions.py
% and inspecting data/sampled/sample_metadata.json.

\section{Appendix C: Expert Annotation Rubric}

The following 1–5 scale was used by the medical co-author (Blessie [Last Name]) to score model reasoning quality. % TODO: Replace [Last Name]

\begin{description}
    \item[1 — Completely wrong reasoning] The model's reasoning demonstrates a fundamental misunderstanding of the medical concept. The selected answer is incorrect and the reasoning, if applied clinically, would be dangerous.
    \item[2 — Mostly wrong] The model identifies some relevant concepts but reaches an incorrect conclusion due to significant errors in reasoning. The answer is wrong and the reasoning shows a substantial gap in medical knowledge.
    \item[3 — Partially correct] The model demonstrates some valid reasoning and may identify the correct differentials or principles, but makes a key error that leads to an incorrect or partially correct conclusion. The answer may be correct for the wrong reason.
    \item[4 — Mostly correct] The model demonstrates sound medical reasoning with only minor gaps or imprecisions. The answer is correct and the reasoning would be appropriate in a clinical education context with minor caveats.
    \item[5 — Excellent] The model's reasoning is clinically appropriate, well-organized, and demonstrates an understanding of the question that would be expected of a well-prepared medical student or physician. The answer is correct and the reasoning adds genuine educational value.
\end{description}

Additionally, annotators recorded:
\begin{itemize}
    \item \textbf{img\_bias\_detected} (yes/no): Whether the model's reasoning reflected assumptions specific to the US healthcare system that would be unfamiliar or misleading for non-US-trained physicians.
    \item \textbf{expert\_notes}: Free-text notes on patterns, errors, or observations not captured by the numeric score.
\end{itemize}

\section{Appendix D: US-Centric Keywords Used for IMG Tagging}

The following keywords were used to automatically tag questions as IMG-relevant. A question was tagged if any keyword appeared in the question stem or any answer option (case-insensitive matching).

\begin{table}[htbp]
\centering
\caption{US-Centric Keyword Lexicon for IMG Tagging}
\label{tab:img_keywords}
\begin{tabular}{ll}
\toprule
Category & Keywords \\
\midrule
Insurance / Financing & medicaid, medicare, insurance, prior authorization, hmo, ppo \\
Referral / Access & primary care physician, referral, emergency medical treatment \\
US Brand-Name Drugs & tylenol, advil, motrin, zofran, ambien, ativan \\
US Guidelines Bodies & acip, uspstf, ama guidelines, joint commission \\
US Demographic Contexts & inner city, rural clinic, community health center \\
\bottomrule
\end{tabular}
\end{table}

This lexicon is conservative by design and likely undercounts IMG-relevant questions. Future work should expand the lexicon through systematic review by physicians trained in non-US medical systems.
